\begin{comment}
\section{Economic Background}

The study of optimal tax policy has deep theoretical foundations in economics, with seminal contributions from \citet{diamond1971optimal}, \citet{mirrlees1976optimal}, and \citet{mankiw2009optimal}. These frameworks seek to design tax systems that optimize social welfare by balancing equity and efficiency considerations. In this work, we situate our analysis within the Mirrleesian framework of non-linear optimal taxation subject to incentive constraints. This approach acknowledges that individuals respond strategically to tax policies, creating a fundamental tension between redistribution and incentives for productivity.

The Mirrleesian framework considers an economy with agents of heterogeneous skills who choose labor supply to maximize their utility. Each agent's utility function increases with consumption (post-tax income) and decreases with labor effort. The government, unable to directly observe skills, must design a tax system based on observable income. This information asymmetry leads to a second-best problem where the optimal tax schedule must respect incentive compatibility constraints.

For analytical tractability, we focus on one-step economies as described by \citet{saez2001using} and \citet{saez2016generalized}. In this simplified setting, we can derive clean expressions for optimal tax rates. Consider an economy with a continuum of agents with skill levels $s \in S$. Each agent chooses labor supply $l$ to earn pre-tax income $z = s \cdot l$. The government implements a tax function $T(z)$, resulting in post-tax income $\hat{z} = z - T(z)$.

\paragraph{Optimal Tax Rate Derivation} Building on \citet{saez2001using}, the optimal marginal tax rate at income level $z$ can be expressed as:
\begin{equation}
T'(z) = \frac{1 - G(z)}{1 - G(z) + a(z)e}
\end{equation}
where $G(z)$ represents the social welfare weight (the value society places on giving an additional dollar to individuals at income level $z$), $a(z)$ is the local Pareto parameter of the income distribution at $z$, and $e$ is the elasticity of taxable income with respect to the net-of-tax rate. For high incomes, the formula simplifies to $\tau^* = \frac{1}{1 + ae}$, where $a$ characterizes the thickness of the upper tail of the income distribution.

Following \citet{arrow1974essays} and \citet{zheng2020ai}, we adopt an isoelastic utility function of the form:
\begin{equation}
u(\hat{z}, l) = \frac{\hat{z}^{1-\eta} - 1}{1-\eta} - \psi(l)
\end{equation}
where $c$ is consumption, $l$ is labor, $\eta$ is the coefficient of relative risk aversion, and $\psi(l)$ represents the disutility of labor. 

\paragraph{Classical Empirical Computation of Saez Optimal Taxes} To empirically determine the Saez optimal tax rates, we follow the approach of \citet{zheng2020ai} and \citet{zheng2022ai}. First, we estimate the income distribution from simulated data, computing the cumulative distribution function $F(z)$ and probability density function $h(z)$. From these, we calculate the Pareto parameter $a(z) = \frac{z \cdot h(z)}{1 - F(z)}$. Second, we estimate the elasticity of taxable income $e$ through a grid search of income responses to tax rate changes across simulation runs. Finally, we compute social welfare weights $G(z)$ based on a social welfare function that encodes society's preferences for redistribution. Plugging these empirical estimates into the Saez formula yields the marginal tax schedule.

\paragraph{Optimal Taxation as Stackelberg Game} The Stackelberg game formulation provides a natural framework for mechanism design in optimal taxation, where the government acts as leader committing to a tax policy while heterogeneous agents respond as followers. This approach aligns with recent advances in computational mechanism design \citep{brero2022learning, brero2022stackelberg}, where the leader's ability to shape follower equilibria through committed strategies proves crucial for welfare optimization. By modeling the tax authority as a Stackelberg leader, we inherit key advantages: the sequential structure \citep{von2010market} allows the policymaker to internalize agents' strategic responses during tax schedule optimization, avoiding circularity problems inherent in simultaneous-move games; this formulation enables principled integration of reinforcement learning methods for equilibrium computation through backward induction analogs; and the Stackelberg paradigm permits modeling bounded rationality via no-regret learning dynamics \citep{brero2022stackelberg}, where followers gradually adapt to the leader's mechanism through experience rather than requiring full rationality assumptions.

% small scenarioes: have stackelberg equilibria
% large: compare the swf in last 10% time steps, average utility distribution (pre and post tax as well)
% saez comparison in flat tax scenario, 

% scenarios: voting, 

% divide swf by pretax income z because otherwise the posttax z component in the social welfare is a constant since the rebate redistributes all tax

\paragraph{LLMs for Persona-Based Utility Estimation} A key limitation in traditional optimal tax models is the assumption of homogeneous utility functions and perfectly rational agents, which can greatly affect designing mechanisms~\citet{rees2018taxing}. As noted by \citet{chetty2009sufficient} and \citet{farhi2010capital}, actual human behavior exhibits various forms of bounded rationality and preference heterogeneity that can significantly impact the welfare implications of tax policies. \citet{hoderlein2004nonparametric} provides a framework for modeling heterogeneous populations, demonstrating that economic demand theory can incorporate diverse preference structures without overly restrictive functional forms. The challenge is particularly pronounced in normative economics, where, as \citet{fleurbaey2004normative} acknowledges, it is standard to "evaluate individual well-being on the basis of self-centered preferences" primarily because it simplifies analysis rather than for its realism. Furthermore,~\citet{maliar2003representative} demonstrate that idiosyncratic uncertainty can have non-negligible effects on aggregate labor-market fluctuations, highlighting the importance of capturing heterogeneity in economic models. Our approach addresses these challenges through the use of LLMs to generate diverse, synthetic human personas with heterogeneous utility functions. This method enables the simulation of human economic behavior (termed "homo silicus"~\citep{horton2023large}) across various scenarios, allowing for the exploration of economic theories through computational simulations with agents endowed with specific preferences and constraints. 

By incorporating synthetic human utility functions, our framework extends the standard utility model to 
\begin{equation*}
    u_{\text{adjusted}} = r \cdot u_{\text{isoelastic}}
\end{equation*}
where $r$ represents a satisfaction factor ranging from 0.5 to 1.0 determined by the LLM based on the agent's persona and current economic situation. As \citet{zhang2009discrete} demonstrates in multigroup growth modeling, the preferences and constraints of heterogeneous households significantly determine wealth and income distribution in competitive economies. Our approach thus enables modeling of the complex interplay between labor choices, tax policies, and individual preferences, particularly for underrepresented populations that traditional economic models often overlook. Through this methodological integration, we advance beyond conventional modeling techniques toward a more realistic representation of economic decision-making in heterogeneous societies.

The use of LLMs in our economic framework addresses a fundamental challenge known as the Lucas Critique \citep{lucas1976econometric}, which argues that policy evaluation based on historical behavioral patterns is invalid when those patterns would change under new policies. LLMs offer a potential solution by providing counterfactual behavioral responses grounded in a broad understanding of human preferences and decision-making. This approach aligns with recent work by \citet{ilut2023economic} on model-based expectations in economic decision-making and \citet{gabaix2020behavioral} on behavioral inattention in economic models.

By modeling the elasticity of taxable income through simulated agent behavior rather than imposing it exogenously, our framework provides a more nuanced view of how different demographic groups might respond to tax changes. This allows us to compute optimal tax schedules that account for the heterogeneity in behavioral responses across the population, potentially leading to more robust and effective tax policies. The LLM-based approach thus represents a significant methodological advance in mechanism design for economic policy, bridging theoretical optimal taxation with empirically grounded behavioral modeling.
\end{comment}

% \section{Optimal Tax Rate Derivation for Piecewise Linear Tax}
% \seth{I converted this directly from Wenzhe's note. Please check if proof is correct.}
% \subsection{Basic Framework}
% Consider a tax bracket \([z_1, z_2]\) with marginal tax rate \(\tau_2\). We analyze three effects of perturbing \(\tau_2 \rightarrow \tau_2 + d\tau\):

% \subsection{Mechanical Tax Increase}
% The immediate revenue gain without behavioral responses:
% \[
% dM = d\tau \cdot (z_2 - z_1) \left[1 - \underbrace{\frac{1}{2}H(z_1) - \frac{1}{2}H(z_2)}_{\text{Population average in bracket}}\right]
% \]
% Where \(H(z)\) is the cumulative distribution function of income.

% \subsection{Social Welfare Effect}
% The welfare loss from increased taxation:
% \[
% dW = -dM \cdot G(z_1,z_2) = -d\tau \cdot (z_2 - z_1)\left[1 - \frac{1}{2}H(z_1) - \frac{1}{2}H(z_2)\right]G(z_1,z_2)
% \]
% With welfare weight:
% \[
% G(z_1,z_2) = \frac{\frac{1}{2}\int_{z_1}^{z_2} g(s)h(s)ds + \int_{z_2}^\infty g(s)h(s)ds}{1 - \frac{1}{2}H(z_1) - \frac{1}{2}H(z_2)}
% \]
% Where \(g(s)\) is the social marginal welfare weight and \(h(s)\) is the income density.

% \subsection{Behavioral Response}
% Taxpayers adjust earnings through elasticity \(e\):
% \[
% \delta z = -d\tau \cdot \frac{e z}{1-\tau_2}
% \]
% Leading to revenue loss:
% \[
% dB = \int_{z_1}^{z_2} \tau_2 \delta z(s) h(s)ds = -d\tau \int_{z_1}^{z_2} \frac{e s \tau_2}{1-\tau_2} h(s)ds
% \]

% \subsection{Optimality Condition}
% Setting net effect to zero (\(dM + dW + dB = 0\)):
% \[
% d\tau \cdot \int_{z_1}^{z_2} \left[(1-\tilde{H})(1 - G) - \frac{e s \tau_2}{1-\tau_2}h(s)\right]ds = 0
% \]
% Where \(\tilde{H}(z_1,z_2) = \frac{1}{2}H(z_1) + \frac{1}{2}H(z_2)\) is the mean CDF.

% \subsection{Key Equation Derivation}
% After simplification:
% \[
% \frac{e\tau_2}{1-\tau_2} \int_{z_1}^{z_2} s h(s)ds = (z_2-z_1)(1-\tilde{H})(1-G)
% \]
% Define Pareto parameter:
% \[
% \alpha(z_1,z_2) = \frac{\int_{z_1}^{z_2} s h(s)ds}{(1-\tilde{H})(z_2-z_1)}
% \]
% Yielding optimal rate:
% \[
% \tau_2 = \frac{1-G}{1-G + \alpha e}
% \]

% \subsection{Continuous Limit Case}
% As bracket width \(dz \to 0\):
% \[
% \begin{aligned}
% \tilde{H} &\to H(z_1) \\
% \alpha &\to \frac{z_1 h(z_1)}{1-H(z_1)} \\
% G &\to \frac{\int_{z_1}^\infty g(s)h(s)ds}{1-H(z_1)}
% \end{aligned}
% \]
% Recovering standard nonlinear tax formula:
% \[
% T'(z) = \frac{1-G(z)}{1-G(z) + \alpha(z)e(z)}
% \]


\section{Background Economic Theory}
In this section, we provide background derivations from Saez economic theory. We note though that the theoretical approach requires strict independence of elasticity between brackets. However, in each bracket, the elasticity parameter depends on the population's behavioral policy, which has a shared dependence across all tax brackets. Additionally, the utility function is assumed to be purely rational. Both of these assumptions are immediately violated in our setting and in real life. These derivations are provided for background and a true analytical solution is not remotely possible. Thus, as noted in our experiments, Saez tax rates require a solution from the LLM Economist to be locally perturbed before finding the optimal policy.

The social welfare:
\begin{align*}
\text{SWF} =& \frac{1}{N} \sum_{i} G_i(u_i(c_i, z_i))\\
c_i(y_i, \bar{r}) =& y_i + \bar{r}\\
y_i(z_i, \tau) =& z_i - T_\tau(z_i)\textbf{}
\end{align*}
Here $N$ is the total number of the agents, $G_i$ is the welfare function, $u_i$ is the utility function, $c_i$ is the post-tax income, $y_i$ is the post-tax income (without rebate), and $z_i$ is the pre-tax income, respectively, for agent $i$. $T_\tau(\cdot)$ is the tax policy parameterized by $\tau$. $\bar{r}$ is the tax rebate for everyone, which can be calculated as:

\begin{equation}
\bar{r}(\tau, z_{1:N}) = \frac{1}{N} \sum_{i} T_\tau(z_i)
\end{equation}

We want to find the tax policy to maximize SWF:
\begin{equation} \label{eq:op_1}
\frac{\dd \text{SWF}}{\dd \tau} = \frac{1}{N} \sum_i G'(u_i) \left[
\frac{\partial u_i}{\partial c_i} \frac{\dd c_i}{\dd \tau} + \frac{\partial u_i}{\partial z_i} \frac{\dd z_i}{\dd \tau}\right] =0
\end{equation}
On the other hand, we know each individual optimizes over $z_i$ to maximize $u_i$, given tax policy $\tau$ (such that ${\dd y_i}/{\dd z_i} = {\partial y_i}/{\partial z_i}$), and rebate $\bar{r}$ (considering $N$ sufficiently large, such that $z_i$ has negligible impact on $\bar{r}$, i.e., $\partial\bar{r}/\partial z_i \approx 0$). This gives:
\begin{equation} \label{eq:op_2}
\frac{\partial u_i}{\partial c_i} \frac{\partial y_i}{\partial z_i} + \frac{\partial u_i}{\partial z_i} = 0
\end{equation}
Finally by chain rule, we have (note that $\tau$ has non-negligible impact on $\bar{r}$):
\begin{equation} \label{eq:op_3}
\frac{\dd c_i}{\dd \tau} = \frac{\partial y_i}{\partial z_i}\frac{\dd z_i}{\dd \tau} +  \frac{\partial y_i}{\partial \tau} + \frac{\dd \bar{r}}{\dd \tau} 
\end{equation}
Define $g_i:= G'(u_i)\cdot\frac{\partial u_i}{\partial c_i} $ and combine \eqref{eq:op_1} \eqref{eq:op_2} \eqref{eq:op_3}, we have: 
\begin{equation*}
    \sum_i g_i \frac{\partial y_i}{\partial \tau} + (\sum_i g_i) \frac{\dd \bar{r}}{\dd \tau}  = 0
\end{equation*}
By rearranging the terms and using the chain rule on the second term, we have:
\begin{equation} \label{eq:general_form}
    \frac{\sum_i g_i \frac{\partial y_i}{\partial \tau}}{\sum_i g_i}  +  \frac{\partial \bar{r}}{\partial \tau}  + \sum_i \frac{\partial \bar{r}}{\partial z_i} \frac{\dd z_i}{\dd \tau} = 0
\end{equation}
or equivalently:
\begin{equation} \label{eq:general_form_difference}
    \underbrace{\frac{\sum_i g_i \frac{\partial y_i}{\partial \tau}}{\sum_i g_i} \dd \tau}_{\dd W} +  \underbrace{\frac{\partial \bar{r}}{\partial \tau} \dd \tau}_{\dd M} + \underbrace{\sum_i \frac{\partial \bar{r}}{\partial z_i} \dd z_i}_{\dd B} = 0
\end{equation}
where $\dd W, \dd M, \dd B$ correspond to social welfare effect, mechanical tax increase, and behavioral response, respectively.

In cases of \eqref{eq:general_form}, it can be simplified to be of the form:
\begin{equation*}
    - A + B - \frac{\tau e}{1-\tau} C = 0
\end{equation*}
Then, we can solve for the tax rate of form:
\begin{equation*}
    \tau = \frac{1-G}{1-G+\alpha\cdot e}
\end{equation*}
where $G = A/B$ and $\alpha = C/B$. We will specify how $G, \alpha, e$ are defined in each case.

\subsection{Top Income Tax Rate~(\citet{saez2001using}, Section 3)}
Consider the tax is \textit{linear} above a given top income $z^*$. Then we can express $y_i$ and $\bar{r}$ as:
\begin{align*}
    y_i(z_i, \tau) &= \begin{cases}
      z_i - T(z_i) & \text{if $z_i < z^*$}\\
      z_i - T(z^*) - \tau(z_i-z^*) & \text{otherwise}
    \end{cases} \\
    \bar{r}(\tau, z_{1:N}) &= \frac{1}{N}\sum_{i:z_i<z^*}T(z_i) + \frac{1}{N}\sum_{i:z_i\geq z^*}\left [T(z_*) + \tau(z_i-z^*)\right ]
\end{align*}
Here $N$ is the total number of the agents, $T(\cdot)$ is the general tax policy for income below $z^*$, and $\tau$ is the \textit{linear} tax rate above $z^*$.

To solve \eqref{eq:general_form}, we first compute the following partial derivatives:
\begin{align*}
    \frac{\partial y_i}{\partial \tau} &= \begin{cases}
      0 & \text{if $z_i < z^*$}\\
      - (z_i-z^*) & \text{otherwise}
    \end{cases} \\
    \frac{\partial \bar{r}}{\partial \tau} &= \frac{1}{N}\sum_{i:z_i\geq z^*}(z_i-z^*) \\
    \frac{\partial \bar{r}}{\partial z_i} &= \frac{1}{N}\tau, \quad \forall i: z_i \geq z^*
\end{align*}
Also we make a key assumption here: \textbf{the change of $\tau$ only affects the income above $z^*$}, i.e.:
\begin{equation*}
    \frac{\dd z_i}{\dd \tau} = 0, \quad \forall i: z_i < z^* 
\end{equation*}
Plug into \eqref{eq:general_form}, we have:
\begin{equation} \label{eq:high_income_form}
    \frac{-\sum_{i:z_i\geq z^*} g_i(z_i-z^*)}{\sum_i g_i}  +  \frac{1}{N}\sum_{i:z_i\geq z^*}(z_i-z^*) + \frac{\tau}{N}\sum_{i:z_i\geq z^*} \frac{\dd z_i}{\dd \tau} = 0
\end{equation}
Assume that the high-income population has the same elasticity:
\begin{equation*}
    e_{z^*} = \frac{(1-\tau)\dd z_i}{z_i\dd(1-\tau)}, \quad \forall i: z_i \geq z^*
\end{equation*}
Then we can rewrite \eqref{eq:high_income_form} as:
\begin{equation*} 
    -\frac{\sum_{i:z_i\geq z^*} g_i(z_i-z^*)}{\sum_i g_i}  +  \frac{1}{N}\sum_{i:z_i\geq z^*} (z_i-z^*) - \frac{\tau e}{1-\tau}\frac{1}{N}\sum_{i:z_i\geq z^*} z_i = 0
\end{equation*}
Therefore we have:
\begin{align*}
    A &= \frac{\sum_{i:z_i\geq z^*} g_i(z_i-z^*)}{\sum_i g_i} \\
    B &= \frac{1}{N}\sum_{i:z_i\geq z^*} (z_i-z^*) \\
    C &= \frac{1}{N}\sum_{i:z_i\geq z^*} z_i
\end{align*}
Such that:
\begin{align*}
    G &= \frac{A}{B} = \frac{\sum_{i:z_i\geq z^*} g_i(z_i-z^*)}{\left [\frac{1}{N}\sum_i g_i\right ] \sum_{i:z_i\geq z^*} (z_i-z^*)} \\
    \alpha &= \frac{C}{B} = \frac{\sum_{i:z_i\geq z^*} z_i}{\sum_{i:z_i\geq z^*} (z_i-z^*)}
\end{align*}
Let $M$ be the number of the agents with income above $z^*$ and define their average income:
\begin{equation*}
    z_M = \frac{1}{M}\sum_{i:z_i\geq z^*} z_i
\end{equation*}
Then we have:
% \seth{this `zd` is confusing. please define $z * d$ where $d$ is the difference operator (is this just log)}
\begin{align*}
    G &= \frac{\frac{1}{M}\sum_{i:z_i\geq z^*} g_i(z_i-z^*)}{\left [\frac{1}{N}\sum_i g_i\right ] (z_M-z^*)} \\
    \alpha &= \frac{z_M}{z_M-z^*} \\
    e_{z^*} &= \frac{(1-\tau)\dd z}{z\dd(1-\tau)}, \quad \forall z \geq z^*
\end{align*}
% \seth{remark if $e_{z^*}$ is tractable}

\paragraph{Remark.} Let $z^*=0$ and $M=N$; thus, for the flat tax case, we can recover the linear income tax rate:
\begin{align*}
    G &= \frac{\sum_{i} g_i z_i}{z_M\sum_i g_i} \\
    \alpha &= 1 \\
    e &= \frac{(1-\tau)\dd z}{z\dd(1-\tau)}, \quad \forall z \geq 0
\end{align*}
As you can see, elasticity is population dependent because $z$ depends on $l$, which depends on the behavioral model of the agents.

\subsection{Non-Linear Income Tax Rate~(\citet{saez2001using}, Section 4)}
For the non-linear tax policy $T(\cdot)$, let $\tau_z=T'(z)$ be the marginal tax rate in $[z,z+\dd z]$. Since it is hard to write down the exact form of $y_i(z_i, \tau)$ and $\bar{r}(\tau, z_{1:N})$ (as $\tau_z$ will affect $T(\cdot)$ in $[z,\infty]$), we use the perturbation approach to compute partial derivatives. 

Consider small $\dd\tau>0$ reform in $[z,z+\dd z]$:
\begin{align*}
    \text{Fix $z_i$:} \quad \dd y_i &= \begin{cases}
      0 & \text{if $z_i < z$}\\
      -(z_i-z)\dd\tau & \text{if $z_i \in [z,z+\dd z]$}\\
      - \dd z\dd\tau & \text{otherwise}
    \end{cases} \\
    \text{Fix $z_{1:N}$:} \quad \dd\bar{r} &= \frac{1}{N}\sum_{i:z_i\geq z}\dd z\dd\tau \\
    \text{Fix $\tau_z$ and $z_{i^-}$:} \quad \dd\bar{r} &= \frac{1}{N}\tau_z\dd z_i, \quad \forall i: z_i \in [z,z+\dd z]
\end{align*}
Therefore, the partial derivatives are:
\begin{align*}
    \frac{\partial y_i}{\partial \tau} &= \begin{cases}
      0 & \text{if $z_i < z$}\\
      -(z_i-z) & \text{if $z_i \in [z,z+\dd z]$}\\
      - \dd z & \text{otherwise}
    \end{cases} \\
    \frac{\partial \bar{r}}{\partial \tau} &= [1-H(z)]\dd z \\
    \frac{\partial \bar{r}}{\partial z_i} &= \frac{1}{N}\tau_z, \quad \forall i: z_i \in [z,z+\dd z]
\end{align*}
Here $H(z)$ is the CDF of $z$.

Also we make a key assumption here: \textbf{the change of $\tau_z$ only affects the income in $[z,z+\dd z]$}, i.e.:
\begin{equation*}
    \frac{\dd z_i}{\dd \tau} = 0, \quad \forall i: z_i \not\in [z,z+\dd z] 
\end{equation*}
Plug into \eqref{eq:general_form}, we have:
\begin{equation} \label{eq:nl_income_form}
    \frac{-\sum_{i:z_i\geq z} g_i\dd z}{\sum_i g_i}  +  [1-H(z)]\dd z + \frac{\tau_z}{N}\sum_{i:z_i \in [z,z+\dd z]} \frac{\dd z_i}{\dd \tau} = 0
\end{equation}

Assume that the population in $[z,z+\dd z]$ has the same elasticity:
\begin{equation*}
    e_z = \frac{(1-\tau_z)\dd z_i}{z_i\dd(1-\tau)}, \quad \forall i: z_i \in [z,z+\dd z]
\end{equation*}
Then we can rewrite \eqref{eq:nl_income_form} as:
\begin{equation*} 
    -\frac{\sum_{i:z_i\geq z} g_i}{\sum_i g_i}  + [1-H(z)] - \frac{\tau_z e_z}{1-\tau_z}\frac{1}{N}\sum_{i:z_i \in [z,z+\dd z]} \frac{z_i}{\dd z} = 0
\end{equation*}
Therefore we have:
\begin{align*}
    A &= \frac{\sum_{i:z_i\geq z} g_i}{\sum_i g_i} \\
    B &= [1-H(z)] \\
    C &= \frac{1}{N}\sum_{i:z_i \in [z,z+\dd z]} \frac{z_i}{\dd z} = zh(z)
\end{align*}
Here $h(z)$ is the PDF of $z$.
Such that:
\begin{align*}
    G &= \frac{A}{B} = \frac{\sum_{i:z_i\geq z} g_i}{\left [\sum_i g_i\right ][1-H(z)]} \\
    \alpha &= \frac{C}{B} = \frac{zh(z)}{[1-H(z)]} \\
    e_z &= \frac{(1-\tau_z)\dd z}{z\dd(1-\tau)}
\end{align*}

% \seth{this requires measuring the elasticity at every point along the tax policy, intractable}

\subsection{Piecewise Linear Income Tax Rate}
For the piecewise linear tax policy $T(\cdot)$, let $\tau_j$ be the marginal tax rate in $j$-th bracket $[z_j,z_{j+1}]$. Following the previous section, we use the perturbation approach to compute partial derivatives. 

Consider small $\dd\tau>0$ reform of $\tau_j$:
\begin{align*}
    \text{Fix $z_i$:} \quad \dd y_i &= \begin{cases}
      0 & \text{if $z_i < z_j$}\\
      -(z_i-z_j)\dd\tau & \text{if $z_i \in [z_j,z_{j+1}]$}\\
      - (z_{j+1}-z_j)\dd\tau & \text{otherwise}
    \end{cases} \\
    \text{Fix $z_{1:N}$:} \quad \dd\bar{r} &= \frac{1}{N}\sum_{i:z_i \in [z_j,z_{j+1}]}(z_i-z_j)\dd\tau + \frac{1}{N}\sum_{i:z_i > z_{j+1}}(z_{j+1}-z_j)\dd\tau \\
    \text{Fix $\tau_j$ and $z_{i^-}$:} \quad \dd\bar{r} &= \frac{1}{N}\tau_j\dd z_i, \quad \forall i: z_i \in [z_j,z_{j+1}]
\end{align*}
Therefore, the partial derivatives are:
\begin{align*}
    \frac{\partial y_i}{\partial \tau} &= \begin{cases}
      0 & \text{if $z_i < z_j$}\\
      -(z_i-z_j) & \text{if $z_i \in [z_j,z_{j+1}]$}\\
      - (z_{j+1}-z_j) & \text{otherwise}
    \end{cases} \\
    \frac{\partial \bar{r}}{\partial \tau} &= \frac{1}{N}\sum_{i:z_i \in [z_j,z_{j+1}]}(z_i-z_j) + \frac{1}{N}\sum_{i:z_i > z_{j+1}}(z_{j+1}-z_j) \\
    \frac{\partial \bar{r}}{\partial z_i} &= \frac{1}{N}\tau_j, \quad \forall i: z_i \in [z_j,z_{j+1}]
\end{align*}
% Here $H(z)$ is the CDF of $z$.

Also we make a key assumption here: \textbf{the change of $\tau_j$ only affects the income in $[z_j,z_{j+1}]$}, i.e.:
\begin{equation*}
    \frac{\dd z_i}{\dd \tau} = 0, \quad \forall i: z_i \not\in [z_j,z_{j+1}] 
\end{equation*}
Plug into \eqref{eq:general_form}, we have:
\begin{align}
    &\frac{-\sum_{i:z_i \in [z_j,z_{j+1}]}g_i(z_i-z_j) - \sum_{i:z_i > z_{j+1}}g_i(z_{j+1}-z_j)}{\sum_i g_i} \notag \\
    &\quad + \frac{1}{N}\left[\sum_{i:z_i \in [z_j,z_{j+1}]}(z_i-z_j) + \sum_{i:z_i > z_{j+1}}(z_{j+1}-z_j)\right] \notag \\
    &\quad + \frac{\tau_j}{N} \sum_{i:z_i \in [z_j,z_{j+1}]} \frac{\dd z_i}{\dd \tau} = 0 \label{eq:pwl_income_form}
\end{align}


Assume that the population in $[z_j,z_{j+1}]$ has the same elasticity:
\begin{equation*}
    e_j = \frac{(1-\tau_j)\dd z_i}{z_i\dd(1-\tau)}, \quad \forall i: z_i \in [z_j,z_{j+1}]
\end{equation*}
Then we can rewrite \eqref{eq:pwl_income_form} as:
\begin{align*} 
    &-\frac{\sum_{i:z_i \in [z_j,z_{j+1}]}g_i(z_i-z_j) + \sum_{i:z_i > z_{j+1}}g_i(z_{j+1}-z_j)}{\sum_i g_i} \\
    &+ \frac{1}{N}\left [\sum_{i:z_i \in [z_j,z_{j+1}]}(z_i-z_j) + \sum_{i:z_i > z_{j+1}}(z_{j+1}-z_j)\right ] \\
    &- \frac{\tau_j e_j}{1-\tau_j}\frac{1}{N}\sum_{i:z_i \in [z_j,z_{j+1}]} z_i \\
    &= 0
\end{align*}
Therefore we have:
\begin{align*}
    A &= \frac{\sum_{i:z_i \in [z_j,z_{j+1}]}g_i(z_i-z_j) + \sum_{i:z_i > z_{j+1}}g_i(z_{j+1}-z_j)}{\sum_i g_i} \\
        &= \frac{\int_{z_j}^{z_{j+1}}h(z)g(z)(z-z_j)\dd z + \int_{z_{j+1}}^{\infty}h(z)g(z)(z_{j+1}-z_j)\dd z}{\int_0^\infty h(z)g(z)\dd z} \\
    B &= \frac{1}{N}\left [\sum_{i:z_i \in [z_j,z_{j+1}]}(z_i-z_j) + \sum_{i:z_i > z_{j+1}}(z_{j+1}-z_j)\right ] \\
        &= [H(z_{j+1})-H(z_j)](z_{M,j}-z_j) + [1-H(z_{j+1})](z_{j+1}-z_j) \\
    C &= \frac{1}{N}\sum_{i:z_i \in [z_j,z_{j+1}]} z_i = \int_{z_j}^{z_{j+1}}h(z)z\dd z
\end{align*}
Here $H(z)$ and $h(z)$ are the CDF and PDF of $z$, and $z_{M,j}$ is the average income in $[z_j,z_{j+1}]$, respectively.
Such that:
\begin{align*}
    G &= \frac{A}{B} = \frac{\int_{z_j}^{z_{j+1}}h(z)g(z)(z-z_j)\dd z + \int_{z_{j+1}}^{\infty}h(z)g(z)(z_{j+1}-z_j)\dd z}{\left [[H(z_{j+1})-H(z_j)](z_{M,j}-z_j) + [1-H(z_{j+1})](z_{j+1}-z_j)\right ]\int_0^\infty h(z)g(z)\dd z} \\
    \alpha &= \frac{C}{B} = \frac{\int_{z_j}^{z_{j+1}}h(z)z\dd z}{[H(z_{j+1})-H(z_j)](z_{M,j}-z_j) + [1-H(z_{j+1})](z_{j+1}-z_j)} \\
    e_{j} &= \frac{(1-\tau_j)\dd z}{z\dd(1-\tau)}, \quad \forall z \in [z_j,z_{j+1}]
\end{align*}

\paragraph{Remark.} Let $z_{j+1}=z_j+\dd z$, we can recover the non-linear income tax rate.
